\documentclass{ximera}

\input{../preamble.tex}


\title[Problem 1]{Problem 1}

\begin{document}
\begin{abstract} \end{abstract}
\maketitle


% Extracted from derivativesAsFunctions.tex, problem #1
\begin{problem}
	
	\begin{enumerate}

    \item
   
      Suppose $f'(2)$ exists. Which of the following must be true?
      \begin{enumerate}
        \item
          $\displaystyle \lim_{x \to 2} f(x)$ must exist, but $\displaystyle \lim_{x \to 2} f(x) \ne f(2)$

        \item
          $\displaystyle \lim_{x \to 2} f(x) = f(2)$.

        \item
          $\displaystyle \lim_{x \to 2} f(x) = f'(2)$

        \item
          $\displaystyle \lim_{x \to 2} f(x)$ need not exist.
      \end{enumerate}
      \begin{explanation}
        The correct answer is (ii):
        \begin{align*}
          \mbox{$f'(2)$ exists} &\iff \mbox{$f$ differentiable at $x = 2$}\\
                                &\implies \mbox{$f$ continuous at $x = 2$}\\
                                &\iff \displaystyle \lim_{x \to 2} f(x) = f(2)
        \end{align*}
      \end{explanation}
\item

      Assuming that $\displaystyle \lim_{x \to 0} \dfrac{\sin(x)}{x} = 1$, which of the following is true?
      \begin{enumerate}
        \item
          $\dfrac{0}{0} = 1$
        \item
          The tangent line to $y = \sin(x)$ at $(0,0)$ has slope $1$.

        \item
          You can cancel the $x$'s.

        \item
          For all $x$ near $0$, $\sin(x) = x$.

        \item
          For all $x$ near $0$, $\sin(x) \approx x$.
      \end{enumerate}
      \begin{explanation}
        This problem has two correct answers: (ii) and (v).

        Statement (ii) is true:
        \begin{align*}
          f'(0) &= \displaystyle \lim_{h \to 0} \dfrac{\sin(0 + h) - \sin (0)}{h} \textrm{ Which has form } \ZeroOverZero \\
                &= \displaystyle \lim_{h \to 0} \dfrac{\sin(h)}{h} \\
                &= 1\\
          &\implies \mbox{slope of tangent line at $(0,0)$ is $1$}
        \end{align*}

        Statement (v) is true:
        When $x$ is near $0$ the tangent line $y = x$ is a good approximation to $f$.
      \end{explanation}
\end{enumerate}

	
\end{problem}



\end{document}
